%% Supplementary material for PMB submission:
%% Reduced viscoelastic FDTD for shear wave propagation
\documentclass[11pt]{article}
\usepackage[a4paper,margin=1in]{geometry}
\usepackage{amsmath,amssymb,bm}
\usepackage{graphicx}
\usepackage{float}
\usepackage{hyperref}
\usepackage{enumitem}
\graphicspath{{./figures/}}

\title{Supplementary Material:\\
A reduced viscoelastic FDTD formulation for ultrasound-driven shear wave
propagation in soft tissue}
\author{Gianmarco F Pinton}
\date{}

\begin{document}
\maketitle

\setcounter{section}{0}
\renewcommand{\thesection}{S\arabic{section}}
\renewcommand{\thesubsection}{S\arabic{section}.\arabic{subsection}}
\renewcommand{\thefigure}{S\arabic{figure}}
\renewcommand{\thetable}{S\arabic{table}}
\renewcommand{\theequation}{S\arabic{equation}}

%% ====================================================================
\section{Theta-Burst TUS Protocol}
\label{sec:supp_theta_burst}

The pipeline supports a theta-burst transcranial ultrasound stimulation (TUS)
protocol motivated by neuromodulation applications.

\paragraph{Protocol parameters.}
The theta-burst envelope consists of 360~$\mu$s tone-bursts repeated at a
1~kHz pulse repetition frequency (PRF), grouped into trains of 3~bursts. Each
3-burst group occupies the first 3~ms of a 200~ms theta period ($1/5$~Hz),
giving an overall duty cycle of ${\sim}$0.5\%. One complete theta cycle (200~ms)
is simulated.

\paragraph{Intensity scaling.}
The acoustic simulation was performed at a source pressure of 1.5~MPa, yielding
a measured focal $I_{\mathrm{SPPA}} \approx 633$~W/cm$^2$ (after CW envelope
correction). Typical neuromodulation protocols operate at lower intensities,
$I_{\mathrm{SPPA}} \approx 20$--$30$~W/cm$^2$. Rather than re-running the
acoustic simulation, the body force is scaled by the ratio of target-to-simulated
focal intensity:
\begin{equation}
  \bm{f}_{\mathrm{neuro}} = \frac{I_{\mathrm{SPPA}}^{\mathrm{target}}}
                                  {I_{\mathrm{SPPA}}^{\mathrm{sim}}}
                            \;\bm{f}_{\mathrm{sim}},
  \label{eq:arf_scaling}
\end{equation}
since the radiation force is proportional to acoustic intensity. For a target of
30~W/cm$^2$ the scaling factor is $30/633 \approx 0.047$, reducing the peak
axial force from ${\sim}$41{,}800 to ${\sim}$2{,}000~N/m$^3$.

\paragraph{Results.}
With the scaled ARF and theta-burst envelope, the shear FDTD solver (12{,}035
time steps over 200~ms) produces a peak displacement of
$|u_x|^{\max} \approx 0.3$~nm at the end of the theta cycle. Instantaneous
displacement during the 3~ms burst window reaches ${\sim}$0.08~$\mu$m before
viscous damping ($\eta = 0.5$~Pa$\cdot$s) dissipates the motion during the
subsequent 197~ms off-period. The sub-nanometer residual at cycle end is
consistent with the very low duty cycle (0.5\%).


%% ====================================================================
\section{From Tissue Displacement to Membrane Tension}
\label{sec:supp_membrane_tension}

In neuromodulation applications, the quantity that ultimately gates
mechanosensitive ion channels is not tissue displacement itself but the
\emph{membrane tension} $T$ (force per unit length, mN/m) experienced by the
lipid bilayer in which the channel is embedded.

The linearized strain tensor
$\varepsilon_{ij} = \tfrac{1}{2}(\partial_i u_j + \partial_j u_i)$ is computed
from the solver displacement field. Because the tissue is nearly
incompressible, the mechanically relevant component is the deviatoric strain
norm $\varepsilon_{\mathrm{dev}}$. The effective membrane tension is then
\begin{equation}
    T = K_A \,\alpha\, \varepsilon_{\mathrm{dev}},
    \label{eq:tension}
\end{equation}
where $K_A$ is the area-expansion modulus of the lipid bilayer (typically
0.1--0.5~N/m for biological membranes) and $\alpha$ is a dimensionless
coupling factor relating far-field tissue strain to local membrane area change.

The tension $T$ enters the Boltzmann gating model for mechanosensitive
channels. Each channel opens with a steady-state probability
$P_o(T) = [1 + \exp(-\Delta A\,(T - T_{1/2})/k_BT_{\mathrm{abs}})]^{-1}$,
where $\Delta A$ is the effective gate area change and $T_{1/2}$ is the
half-activation tension.

For the channels most relevant to ultrasonic neuromodulation, the
best-characterized half-activation tensions are: Piezo1,
$T_{1/2} = 2.7 \pm 0.4$~mN/m; and TRAAK, $T_{1/2} \approx 5$~mN/m. Values
for TREK-1 (${\sim}$3.5~mN/m) and TREK-2 (${\sim}$4~mN/m) carry larger
uncertainty.

These tension scales set the target for the simulation pipeline: the shear wave
solver must resolve displacement gradients accurately enough that the resulting
membrane tensions can be compared with the channel $T_{1/2}$ values. For a
representative $K_A = 0.25$~N/m and $\alpha = 1$, a deviatoric strain of
${\sim}10^{-5}$ produces a membrane tension of ${\sim}2.5$~$\mu$N/m---well
below channel thresholds---while strains of order $10^{-2}$ reach the
few-mN/m range where channel gating becomes significant.


\bibliographystyle{unsrt}
\bibliography{general}

\end{document}
